% !TeX root = part6.tex
\documentclass[../final.tex]{subfiles}

\begin{document}

\practice{6}{}{Термы многоэлектронного атома}

\begingroup
\renewcommand{\thesection}{4.\arabic{section}}

\rtheory{LS-связь, принцип Паули и правила Хунда}

В приближении \rconcept{LS-связи} сначала складываем орбитальные
моменты электронов и отдельно их спины:
\[
    \mathbf L=\sum_i\mathbf l_i,\qquad
    \mathbf S=\sum_i\mathbf s_i,\qquad
    \mathbf J=\mathbf L+\mathbf S.
\]
Терм обозначается ${}^{2S+1}L_J$. Буквам $S,P,D,F,G,H$
соответствуют $L=0,1,2,3,4,5$. Для заданных $L,S$
\[
    J=|L-S|,\ |L-S|+1,\ldots,L+S.
\]
Без учёта тонкой структуры число состояний терма равно
\[
    g_{LS}=(2L+1)(2S+1)=\sum_J(2J+1).
\]
Замкнутая подоболочка имеет $L=S=0$. Чётность конфигурации равна
$\pi=(-1)^{\sum_i l_i}$; нечётную чётность обозначаем знаком $\mathrm{o}$.

\rconcept{Эквивалентные электроны} имеют одинаковые $n,l$.
Для них недостаточно формально сложить моменты: необходимо исключить
состояния, запрещённые принципом Паули. Число независимых
антисимметризованных состояний подоболочки $nl^N$ равно
\[
    \mathcal N=\binom{2(2l+1)}{N}.
\]
Электроны разных подоболочек также подчиняются принципу Паули,
но при сложении уже допустимых термов этих подоболочек новых
ограничений на $L,S$ не возникает.

\subsection*{Как выбирать старший терм при подсчёте}
Составляем таблицу числа состояний с заданными
$M_L=\sum m_{l_i}$ и $M_S=\sum m_{s_i}$.
Выбираем наибольшее оставшееся $M_S$, затем наибольшее $M_L$
при этом $M_S$. Получаем $S=M_S$, $L=M_L$ и вычитаем по одному
состоянию из каждой клетки прямоугольника
$-S\leq M_S\leq S$, $-L\leq M_L\leq L$. Повторяем до опустошения таблицы.
Вычитаем \rconcept{числа состояний}: функция с определёнными $L,S$
обычно является линейной комбинацией детерминантов.

\subsection*{Как выбирать основной терм}
Правила Хунда применяем последовательно: максимальный $S$;
при этом $S$ максимальный $L$; затем минимальный $J=|L-S|$
для подоболочки, заполненной менее чем наполовину, и максимальный
$J=L+S$ для подоболочки, заполненной более чем наполовину.
Для наполовину заполненной подоболочки основной терм имеет $L=0$.
Старший терм при подсчёте не обязательно имеет наименьшую энергию.

\clearpage
\task[1]{Термы конфигурации \texorpdfstring{$np^3$}{np3}}
\label{task:fiham6_p3}

\condition
Найти все термы атома, единственная незаполненная электронная
подоболочка которого имеет конфигурацию $np^3$.

\solution
Для $p$-электрона $l=1$, $m_l=1,0,-1$, $m_s=\pm1/2$.
Введём шесть спин-орбиталей:
\[
\begin{array}{c|rrrrrr}
    &A&B&C&A'&B'&C'\\ \hline
    m_l&1&0&-1&1&0&-1\\
    m_s&1/2&1/2&1/2&-1/2&-1/2&-1/2
\end{array}
\]
Запись $ABC$ далее означает нормированный слэтеровский детерминант,
а не обычное произведение одноэлектронных функций. Перестановка
его столбцов меняет только знак, поэтому $ABC$ и $CBA$
не являются разными состояниями. Всего
\[
    \mathcal N=\binom63=20.
\]

\subsection*{Таблица микросостояний}
При $M_S=3/2$ все три спина направлены вверх. Принцип Паули
оставляет только $ABC$ с $M_L=0$. Аналогично при $M_S=-3/2$
есть только $A'B'C'$ с $M_L=0$.
При $M_S=1/2$ два электрона имеют спин вверх, один вниз:
\begin{center}
\begin{tabular}{c|l|c}
    \rconcept{$M_L$}&\rconcept{Детерминанты при $M_S=1/2$}&
        \rconcept{Число}\\ \hline
    $2$&$ABA'$&1\\
    $1$&$ABB',\ ACA'$&2\\
    $0$&$ABC',\ ACB',\ BCA'$&3\\
    $-1$&$ACC',\ BCB'$&2\\
    $-2$&$BCC'$&1
\end{tabular}
\end{center}
Смена направлений всех спинов даёт ещё девять состояний с $M_S=-1/2$.
Полная таблица чисел состояний имеет вид
\begin{center}
\begin{tabular}{c|rrrrr}
    $M_S\backslash M_L$&$-2$&$-1$&$0$&$1$&$2$\\ \hline
    $3/2$ &0&0&1&0&0\\
    $1/2$ &1&2&3&2&1\\
    $-1/2$&1&2&3&2&1\\
    $-3/2$&0&0&1&0&0
\end{tabular}
\end{center}

\clearpage
\subsection*{Выделение термов}
Состояние $(M_L,M_S)=(0,3/2)$ порождает $L=0,S=3/2$, то есть
${}^4S$. Вычитаем его четыре состояния: по одному при $M_L=0$
и $M_S=3/2,1/2,-1/2,-3/2$.

В оставшейся таблице максимальные проекции равны
$M_S=1/2$, $M_L=2$. Это терм ${}^2D$ с десятью состояниями.
Вычитаем по одному состоянию при каждом $M_L=-2,-1,0,1,2$
в обеих строках $M_S=\pm1/2$.
Остаётся по одному состоянию при $M_L=-1,0,1$, $M_S=\pm1/2$:
это ${}^2P$ с шестью состояниями.
\[
    4+10+6=20=\binom63.
\]
Все три электрона имеют нечётное $l=1$, поэтому термы нечётные.

\answer
\[
    np^3:\quad {}^4S^{\mathrm o},\quad{}^2D^{\mathrm o},\quad{}^2P^{\mathrm o}.
\]
С учётом тонкой структуры получаем пять уровней:
\[
    {}^4S^{\mathrm o}_{3/2},\qquad
    {}^2D^{\mathrm o}_{3/2,5/2},\qquad
    {}^2P^{\mathrm o}_{1/2,3/2}.
\]
По первому правилу Хунда основной терм --- ${}^4S^{\mathrm o}_{3/2}$.

\clearpage
\task[2]{Термы двух конфигураций с неэквивалентными электронами}

\condition
Выписать возможные типы термов конфигураций
$np^1n'p^2$ и $ns^1n'd^2$. В первой конфигурации $n\ne n'$.

\solution
\subsection*{Конфигурация \texorpdfstring{$np^1n'p^2$}{np1 n'p2}}
Сначала находим допустимые родительские термы пары $n'p^2$.
При перестановке двух одинаковых орбитальных моментов $l$
пространственная функция с заданным $L$ приобретает множитель
$(-1)^{2l-L}=(-1)^L$. При чётном $L$ нужна антисимметричная
спиновая функция ($S=0$), при нечётном --- симметричная ($S=1$).
Для двух $p$-электронов $L=0,1,2$, поэтому
\[
    n'p^2:\quad {}^1S,\ {}^1D,\ {}^3P.
\]
Проверка: $1+5+9=15=\binom62$.
Один $np$-электрон имеет $l=1,s=1/2$, то есть терм ${}^2P$.
Складываем отдельно $L$ и $S$ каждого родительского терма
с моментами этого электрона:
\begin{center}
\begin{tabular}{c|c|c|l}
    \rconcept{Родитель $n'p^2$}&$L$&$S$&\rconcept{Дочерние термы}\\ \hline
    ${}^1S$&1&$1/2$&${}^2P$\\
    ${}^1D$&1,2,3&$1/2$&${}^2P,\ {}^2D,\ {}^2F$\\
    ${}^3P$&0,1,2&$1/2,3/2$&
        ${}^2S,\ {}^2P,\ {}^2D,\ {}^4S,\ {}^4P,\ {}^4D$
\end{tabular}
\end{center}
Итак, типов термов семь. Однако ${}^2P$ возникает из трёх разных
родителей, а ${}^2D$ из двух. Это независимые термы одинакового
типа; межэлектронное взаимодействие может смешивать их при
одинаковых $L,S$. Всего с учётом повторений десять LS-термов.
Все термы конфигурации нечётные.

\answer
\[
    np^1n'p^2:\quad
    {}^2S^{\mathrm o},\ {}^2P^{\mathrm o},\ {}^2D^{\mathrm o},\ {}^2F^{\mathrm o},\
    {}^4S^{\mathrm o},\ {}^4P^{\mathrm o},\ {}^4D^{\mathrm o}.
\]
Возможные значения $J$ и число повторений:
\begin{center}
\begin{tabular}{c|c|c|r}
    \rconcept{Тип}&\rconcept{Число термов}&$J$&
        \rconcept{Всего состояний}\\ \hline
    ${}^2S^{\mathrm o}$&1&$1/2$&2\\
    ${}^2P^{\mathrm o}$&3&$1/2,3/2$&18\\
    ${}^2D^{\mathrm o}$&2&$3/2,5/2$&20\\
    ${}^2F^{\mathrm o}$&1&$5/2,7/2$&14\\
    ${}^4S^{\mathrm o}$&1&$3/2$&4\\
    ${}^4P^{\mathrm o}$&1&$1/2,3/2,5/2$&12\\
    ${}^4D^{\mathrm o}$&1&$1/2,3/2,5/2,7/2$&20
\end{tabular}
\end{center}
Проверка полноты:
\[
    2+18+20+14+4+12+20=90=6\binom62.
\]
При $n=n'$ получилась бы конфигурация $np^3$, и ответ следовало
бы брать из задачи~\ref{task:fiham6_p3}, а не из этой таблицы.

\subsection*{Конфигурация \texorpdfstring{$ns^1n'd^2$}{ns1 n'd2}}
Для пары эквивалентных $d$-электронов $l=2$, $L=0,1,2,3,4$.
Та же проверка симметрии даёт
\[
    n'd^2:\quad {}^1S,\ {}^1D,\ {}^1G,\ {}^3P,\ {}^3F.
\]
Число состояний пары равно
\[
    1+5+9+9+21=45=\binom{10}2.
\]
Для $ns$-электрона $l=0$, $s=1/2$. Поэтому итоговый $L$ равен
орбитальному моменту родителя, синглет даёт только дублет,
а триплет --- дублет и квартет:
\begin{center}
\begin{tabular}{c|l|c}
    \rconcept{Родитель $n'd^2$}&\rconcept{Дочерние термы}&$J$\\ \hline
    ${}^1S$&${}^2S$&$1/2$\\
    ${}^1D$&${}^2D$&$3/2,5/2$\\
    ${}^1G$&${}^2G$&$7/2,9/2$\\
    ${}^3P$&${}^2P;\ {}^4P$&$1/2,3/2;\quad1/2,3/2,5/2$\\
    ${}^3F$&${}^2F;\ {}^4F$&$5/2,7/2;\quad3/2,5/2,7/2,9/2$
\end{tabular}
\end{center}

\answer
\[
    ns^1n'd^2:\quad {}^2S,\ {}^2P,\ {}^2D,\ {}^2F,\ {}^2G,\ {}^4P,\ {}^4F.
\]
Все семь термов чётные; каждый встречается один раз.
Проверка:
\[
    2+6+10+14+18+12+28=90=2\binom{10}2.
\]

\clearpage
\task[3]{Число электронов по основному терму}

\condition
Определить число электронов в единственной незаполненной
подоболочке атома, если основной терм равен:
а) ${}^3F_2$; б) ${}^6S_{5/2}$.

\solution
Пусть орбитальный момент подоболочки $l$, число электронов $N$,
а вместимость $q=2(2l+1)$. По первому правилу Хунда число
неспаренных электронов в основном состоянии равно
\[
    k=\min(N,q-N)=2S.
\]
Можно описывать либо $k$ электронов в менее чем наполовину
заполненной подоболочке, либо $k$ дырок в более чем наполовину
заполненной. Максимальный орбитальный момент при одинаковых
спинах получаем, занимая различные $m_l=l,l-1,\ldots,l-k+1$:
\[
    L_{\max}=kl-\frac{k(k-1)}2.
\]

\subsection*{а) Основной терм \texorpdfstring{${}^3F_2$}{3F2}}
Мультиплетность $3$ означает $S=1$, следовательно $k=2$.
Буква $F$ означает $L=3$. По второму правилу Хунда
\[
    3=2l-1\quad\Rightarrow\quad l=2.
\]
Это $d$-подоболочка с вместимостью десять электронов.
Возможны $d^2$ и $d^8$ (два электрона или две дырки), но
третий закон Хунда различает их:
\[
    d^2:\quad J_{\text{осн}}=L-S=2,\qquad
    d^8:\quad J_{\text{осн}}=L+S=4.
\]
Указано $J=2$, поэтому подходит только $d^2$.

\subsection*{б) Основной терм \texorpdfstring{${}^6S_{5/2}$}{6S5/2}}
Мультиплетность $6$ означает $S=5/2$, следовательно $k=5$.
Для $S$-терма $L=0$. Поэтому
\[
    0=5l-\frac{5\cdot4}2=5l-10\quad\Rightarrow\quad l=2.
\]
Подоболочка вновь $d$, но теперь наполовину заполненная:
$N=5=q-N$. Пять параллельных спинов занимают все пять
орбитальных состояний, и
\[
    M_L=2+1+0-1-2=0,\qquad S=5/2,\qquad J=S=5/2.
\]

\answer
\[
    {}^3F_2:\quad nd^2,\ N=2;\qquad
    {}^6S_{5/2}:\quad nd^5,\ N=5.
\]
Главное квантовое число $n$ из обозначения терма определить нельзя.
Вывод использует условие о \rassumption{единственной незаполненной
подоболочке} и правила Хунда для основного состояния.

\clearpage
\task[4]{Спин ядра \texorpdfstring{${}^{59}\mathrm{Co}$}{59Co}}

\condition
Определить спин ядра ${}^{59}\mathrm{Co}$, если основной
электронный уровень атома содержит восемь компонент
сверхтонкого расщепления (уровней с различными $F$).

\solution
\subsection*{Электронный момент основного уровня}
Основная конфигурация нейтрального кобальта
\[
    \mathrm{Co}:\quad [\mathrm{Ar}]\,3d^7 4s^2.
\]
Замкнутые подоболочки имеют $L=S=0$. У $3d^7$ три дырки,
поэтому $S=3/2$. Для максимального $M_L$ пять спинов вверх занимают
все $m_l$, а два спина вниз --- $m_l=2,1$:
\[
    M_L=(-2-1+0+1+2)+(2+1)=3,\qquad L=3.
\]
Для терма ${}^4F$ подоболочка заполнена более чем наполовину:
\[
    J=L+S=3+\frac32=\frac92,\qquad
    \text{основной уровень }{}^4F_{9/2}.
\]

\subsection*{Число сверхтонких уровней}
При наличии спина ядра $I$ полный момент атома равен
$\mathbf F=\mathbf J+\mathbf I$, а его значения
\[
    F=|J-I|,\ |J-I|+1,\ldots,J+I.
\]
Число различных значений $F$ составляет
\[
    N_F=J+I-|J-I|+1=2\min(I,J)+1.
\]
Если бы $I\geq J=9/2$, получили бы $N_F=2J+1=10$,
что не согласуется с условием. Значит, $I<J$ и
\[
    8=2I+1\quad\Rightarrow\quad I=\frac72.
\]
Полученные уровни имеют
\[
    F=1,2,3,4,5,6,7,8.
\]
Проверка числа магнитных состояний:
$\sum_{F=1}^{8}(2F+1)=80=(2I+1)(2J+1)$.

\answer
\[
    I({}^{59}\mathrm{Co})=\frac72.
\]
\endgroup
\end{document}
